\documentclass[11pt,letterpaper]{article}
\usepackage[margin=1in]{geometry}
\usepackage[T1]{fontenc}
\usepackage{amsmath,amssymb}
\usepackage{enumitem}
\usepackage{titlesec}
\usepackage{xcolor}
\usepackage{array}
\usepackage[colorlinks=false]{hyperref}
 
\definecolor{hd}{HTML}{1F3864}
\titleformat{\section}{\large\bfseries\color{hd}}{}{0pt}{}
\titlespacing*{\section}{0pt}{14pt}{4pt}
\setlist[itemize]{leftmargin=1.4em,itemsep=2pt,topsep=3pt}
\setlength{\parindent}{0pt}
\setlength{\parskip}{5pt}
 
\newcommand{\curl}{\nabla\times}
\newcommand{\dv}{\nabla\cdot}
\renewcommand{\vec}[1]{\mathbf{#1}}
 
\begin{document}
 
{\LARGE\bfseries Exam 3 Outline}\\[2pt]
\textit{(For instructors only, to help lead the review. Do not share with students.)}
 
\textbf{Total: 100 points}
 
%------------------------------------------------------------
\section*{1. Calculate $\curl\vec F$ and $\dv\vec F$, and determine whether $\vec F$ is conservative (10 points)}
\begin{itemize}
  \item 4 pts: each calculation ($\curl\vec F$ and $\dv\vec F$).
  \item 2 pts: correctly determining whether $\vec F$ is conservative.
\end{itemize}
 
\section*{2. Use Green's Theorem to compute a line integral (12 points)}
\begin{itemize}
  \item 4 pts: Calculate the $\mathbf{k}$-component of $\curl\vec F$.
  \item 4 pts: Set up the double integral correctly (bounds of integration).
  \item 4 pts: Evaluate the double integral correctly.
  \item Note: No credit if Green's Theorem is stated incorrectly or if the line integral is computed by direct parameterization.
  \item Review reminder: Have students practice evaluating double integrals in both Cartesian and polar coordinates. Students often got stuck on this step in previous semesters.
\end{itemize}
 
\section*{3. Find a parameterization for each surface (12 points)}
\textit{(Similar to HW9, Problem 2)}
 
Three surfaces are given. Each parameterization is worth 4 points:
\begin{itemize}
  \item 3 pts: Correct parameterization $\vec r(u,v)$ of the surface.
  \item 1 pt: Correct bounds for each parameter.
  \item No credit for a parameterization that uses anything other than two parameters.
  \item Review reminders: (1) Review parameterizations of common surfaces: cone, cylinder, sphere, plane, and surface of revolution. (2) A surface parameterization has exactly two parameters, and bounds must be given for each.
\end{itemize}
 
\section*{4. Set up, but do not evaluate, a surface integral of a scalar function (12 points)}
The parameterization $\vec r(u,v)$ is given.
\begin{itemize}
  \item \textbf{5 pts: Overall procedure} (presence and sequence, not correctness):
  \begin{itemize}
    \item 1 pt: Find $\vec r_u$ and $\vec r_v$.
    \item 1 pt: Find $\vec r_u\times\vec r_v$.
    \item 1 pt: Find the magnitude $|\vec r_u\times\vec r_v|$.
    \item 1 pt: Substitute $\vec r(u,v)$ into $f(x,y,z)$.
    \item 1 pt: Give bounds for $u$ and $v$ and write the double integral.
  \end{itemize}
  \item \textbf{7 pts: Accuracy of each step} (graded independently against the key):
  \begin{itemize}
    \item 2 pts: Find $\vec r_u\times\vec r_v$ correctly (including $\vec r_u$ and $\vec r_v$).
    \item 1 pt: Find $|\vec r_u\times\vec r_v|$ correctly.
    \item 2 pts: Substitute $\vec r(u,v)$ into $f$ correctly.
    \item 2 pts: Give correct bounds for $u$ and $v$ and write the correct double integral (including the factor $|\vec r_u\times\vec r_v|$ and $du\,dv$).
  \end{itemize}
  \item \textbf{Policy:} See the Grader Notes. Algebraic-error cap: 1 pt per independent slip, maximum deduction 3 pts, minimum score 9/12. If any conceptual or setup error is present (e.g., wrong bounds, omitting $|\vec r_u\times\vec r_v|$, or integrating a vector expression instead of a scalar), the cap does not apply and each step is graded against the key.
\end{itemize}
 
\section*{5. Compute a flux integral of a vector field: set up and evaluate (20 points)}
\begin{itemize}
  \item \textbf{8 pts: Overall procedure} (presence and sequence, not correctness):
  \begin{itemize}
    \item 1 pt: Set up $\vec r(u,v)$ for the oriented surface.
    \item 2 pts: Find $\vec r_u\times\vec r_v$ and check orientation.
    \item 1 pt: Substitute $\vec r(u,v)$ into $\vec F$.
    \item 1 pt: Take the dot product.
    \item 1 pt: Give bounds for $u$ and $v$.
    \item 2 pts: Set up and attempt the double integral.
  \end{itemize}
  \item \textbf{12 pts: Accuracy of each step} (graded independently against the key):
  \begin{itemize}
    \item 2 pts: Set up $\vec r(u,v)$ correctly.
    \item 4 pts: Find $\vec r_u\times\vec r_v$ correctly and check orientation ($-2$ if orientation is not checked).
    \item 3 pts: Dot product and bounds:
    \begin{itemize}
      \item 2 pts: Compute $\vec F(\vec r(u,v))\cdot(\vec r_u\times\vec r_v)$ correctly.
      \item 1 pt: Give correct bounds for $u$ and $v$.
    \end{itemize}
    \item 3 pts: Evaluate the double integral correctly.
  \end{itemize}
  \item \textbf{Policy:} See the Grader Notes. Algebraic-error cap: 2 pts per independent slip, maximum deduction 6 pts, minimum score 14/20. If any conceptual or setup error is present, the cap does not apply and each step is graded against the key. Partial credit for method comes from the procedure points.
  \item \textit{Grader examples:}
  \begin{itemize}
    \item A correct $\vec r$ and orientation check, but a sign error in $\vec r_u\times\vec r_v$ carried through: one slip, so 18/20.
    \item A sign error in the dot product plus an arithmetic slip in the final integral: two independent slips, so 16/20.
  \end{itemize}
\end{itemize}
 \newpage
\section*{6. Use Stokes' Theorem to evaluate $\displaystyle\iint_S (\curl\vec F)\cdot d\vec S$ (20 points)}
\begin{itemize}
  \item No credit if Stokes' Theorem is not used.
  \item \textbf{8 pts: Overall procedure} (presence and sequence, not correctness):
  \begin{itemize}
    \item 1 pt: Identify the boundary curve.
    \item 1 pt: Parameterize the curve with one parameter.
    \item 1 pt: Find $\vec r\,'(t)$.
    \item 1 pt: Substitute $\vec r(t)$ into $\vec F$.
    \item 1 pt: Take the dot product.
    \item 1 pt: Give bounds for $t$.
    \item 2 pts: Set up and attempt the single-variable integral.
    \item If $\vec r(t)$ uses more than one parameter, cap procedure credit at 4 pts.
  \end{itemize}
  \item \textbf{12 pts: Accuracy of each step} (graded independently against the key):
  \begin{itemize}
    \item 3 pts: Identify the boundary curve and give a one-parameter $\vec r(t)$ correctly.
    \item 2 pts: Find $\vec r\,'(t)$ correctly.
    \item 2 pts: Substitute $\vec r(t)$ into $\vec F$ correctly.
    \item 2 pts: Compute $\vec F(\vec r(t))\cdot\vec r\,'(t)$ correctly.
    \item 3 pts: Evaluate the single-variable integral correctly (including correct bounds for $t$).
  \end{itemize}
  \item \textbf{Policy:} Same as Problem 5 (see the Grader Notes). Algebraic-error cap: 2 pts per independent slip, maximum deduction 6 pts, minimum score 14/20. Wrong bounds for $t$ is a setup error, so the cap does not apply in that case.
  \item Review reminder: Tell students that solving $\iint_S(\curl\vec F)\cdot d\vec S$ by direct parameterization of the surface, rather than by Stokes' Theorem, earns no credit.
\end{itemize}
 
\section*{7. Conceptual questions ($\sim$14 points)}
Questions are to be determined; suggestions welcome. Possible topics:
\begin{itemize}
  \item Choosing among the Fundamental Theorem of Line Integrals, Green's Theorem, Stokes' Theorem, and direct parameterization: select all methods that could be used to evaluate a given line integral.
  \item The physical meaning of curl (rotation) and divergence (source or sink).
  \item Identifying which expressions involving the gradient, curl, and divergence are well-defined: select all that apply (similar to HW9, Problem 19).
\end{itemize}
Practice problems for the conceptual questions will be added to the current exam practice materials, which will be sent out today.
 
%------------------------------------------------------------
\clearpage
\section*{Grader Notes (Problems 4, 5, and 6)}
 
\textbf{Grading policy.} Grade each accuracy step independently against the answer key. Grading does not stop after an error, but a later step that is correct only relative to the student's own earlier error earns no accuracy credit. Procedure points reward a real attempt at each step in the correct order, so grade them for presence and sequence, not correctness. Do not award a step's procedure points if the work shown is nonsensical or unrelated to the problem.
 
\textbf{Algebraic-error cap.} If all errors are algebraic or arithmetic (e.g., a sign error, a dropped factor, a miscomputed derivative, a wrong simplification, or an arithmetic slip in evaluating the integral), deduct a fixed amount per independent slip, up to a maximum deduction. Do not penalize the same slip again in later steps.
 
\begin{center}
\renewcommand{\arraystretch}{1.3}
\begin{tabular}{|c|c|c|c|c|}
\hline
\textbf{Problem} & \textbf{Points} & \textbf{Deduction per slip} & \textbf{Max deduction} & \textbf{Minimum score}\\
\hline
4 & 12 & 1 pt & 3 pts & 9/12\\
5 & 20 & 2 pts & 6 pts & 14/20\\
6 & 20 & 2 pts & 6 pts & 14/20\\
\hline
\end{tabular}
\end{center}
 
\textbf{The cap does not apply} if there is any conceptual or setup error, such as a wrong $\vec r(u,v)$ or $\vec r(t)$, the wrong number of parameters, wrong bounds, a missing or incorrect orientation check, a wrong formula, or not using the required theorem. In that case, grade each step against the key with no cap.
 
\end{document}